% ===================================================================
%  పట్టిక సంఖ్య 12 — స్టేషను. రైలుమార్గాలు. ఓడలు.
%
%  Traduction, faite en 2026, en telougou d'aujourd'hui, sur l'ido de
%  texto/io. Regles de la colonne : voir 10-pattika-01.tex. Une seule
%  suite d'alineas ; le « Noto. » d'ouverture appartient au bloc apar,
%  et la note de bas de page du feuillet 85 a son bloc noto, aux
%  memes places qu'en hindi.
%
%  LE RENVOI (*) N'EST PAS UN NUMERO ET SE PORTE QUAND MEME. Le bloc
%  19 appelle la note par « \textsuperscript{(*)} » ; renvoji.py le
%  releve comme les autres, et la colonne l'ecrit tel quel.
%
%  LE TELOUGOU DU RAIL EST ENTIER, ET IL EST ANCIEN. Le chemin de fer
%  est arrive en Andhra en 1893, trente ans avant ce livret : la
%  langue a eu le temps de se faire ses mots, et elle ne les a pas
%  empruntes. పట్టాలు les rails (64), బోగీ le wagon (37), కూత le
%  sifflet (59), సొరంగం le tunnel (87), పట్టాల చీలిక l'aiguillage
%  (78), సామాను le bagage (17), మోతకూలీ le porteur (16), వేళల పట్టిక
%  l'indicateur (26). C'est le deuxieme tableau de suite ou rien ne
%  manque, apres le port — et pour la meme raison : la chose est la,
%  donc le mot est la. Seuls restent en anglais les objets que le
%  telougou nomme en anglais aujourd'hui : ప్లాట్‌ఫారం, ఇంజను,
%  బాయిలరు, పిస్టన్లు, సూట్‌కేసు.
%
%  TREIZE RETOURNEMENTS SUR LES VINGT-HUIT PAIRES DE parigi.py.
%  Rien de nouveau dans le tri des quinze autres : enumeration,
%  renvoi de gauche deja modificateur, frontiere de phrase.
%
%  ET UN FAUX POSITIF QUI A APPRIS QUELQUE CHOSE A L'OUTIL. La liste
%  donnait « (3) ← jonti ← (1) » au bloc 02-1 : « jonti » n'est pas un
%  mot, c'est la queue de « voyajonti », les voyageurs. En cherchant
%  d'ou venait ce debris on a trouve mieux qu'un cas : parigi.py
%  depouillait « \cc » comme n'importe quelle macro, en le remplacant
%  par une espace. Or « \cc » n'est pas une coupure, c'est une
%  SOUDURE — le fac-simile coupe ses mots en fin de ligne et la
%  transcription le note ainsi. « ku\cc ranta » etait donc lu
%  « ku ranta », et SIX paires du livret se trouvaient tenues par des
%  fragments qui ressemblaient a des participes : ranta, sante, mante,
%  vinta, dante, danta. Elles etaient justes par accident. La soudure
%  faite avant le depouillement, ce sont maintenant ekiranta,
%  impulsante, arivinta, fumante, ludante, sidanta qui les tiennent —
%  les vrais mots. Le livret passe de 385 a 386 paires.
%
%  « voyajonti », lui, reste signale, et c'est irreductible : en ido
%  le nom d'agent pluriel et l'adjectif participe ont la meme finale
%  (-anti, -onti). Aucune expression reguliere ne les separe. On le
%  laisse donc, et on l'ecrit ici plutot que d'elargir le motif pour
%  le faire taire : c'est du bruit connu, ce qui n'est pas la meme
%  chose que du bruit ignore.
% ===================================================================

%%K t12-apar-1 apar
\VUsaut{4.0mm}
\VUtitre{15.0pt}{60}{పట్టిక సంఖ్య 12}
\VUsaut{2.4mm}
\VUpk{11.4pt}{40}{స్టేషను. --- రైలుమార్గాలు. --- ఓడలు.}
\VUsaut{3.4mm}
\textsc{గమనిక.} --- \textit{ప్రతి దేశంలోనూ వేళల లెక్క వేరువేరుగా ఉంటుంది; అందుకే
ఉదాహరణగా మేము} \VUgras{ఫ్రెంచి పట్టిక} \textit{వేళలను తీసుకున్నాము.}

%%K t12-01-1 p
1. --- నేను ఇప్పుడే జర్మనీ గుండా ప్రయాణించాను. బయలుదేరే ముందు, రైళ్ళు ఎప్పుడు
బయలుదేరతాయో తెలుసుకునేందుకు \VUgras{రైలు వేళల పుస్తకం} \textsuperscript{(15)} కొన్నాను.
అన్నిటికంటే అనుకూలమైన రైలు పారిస్‌నుంచి రాత్రి తొమ్మిది గంటలకు బయలుదేరి, స్ట్రాస్‌బూర్గ్‌కు
ఒంటిగంట పదమూడు నిమిషాలకు చేరుతుందని తేల్చుకున్నాక, నా ప్రయాణం సిద్ధం చేసుకున్నాను;
మరుసటి రోజు స్టేషనుకు తీసుకెళ్ళమని చెప్పాను.

%%K t12-02-1 p
2. --- పెద్ద బయలుదేరు హాలులోకి నేను వెళ్ళేసరికి, పల్లె \VUgras{మతగురువు}
\textsuperscript{(2)} వెనుక — ఆయన తన \VUgras{ప్రయాణపు సంచిని} \textsuperscript{(3)}
చేతిలో పట్టుకున్నారు — చాలామంది \VUgras{ప్రయాణికులు} \textsuperscript{(1)} అప్పటికే
వచ్చేశారు.

%%K t12-02-2 p
నేను \VUgras{తంతి (టెలిగ్రామ్)} \textsuperscript{(5)} పంపాలనుకున్నాను; అందుకని
\VUgras{తనిఖీదారుడిని} \textsuperscript{(6)} అడిగి \VUgras{తంతి కార్యాలయం}
\textsuperscript{(4)} ఎక్కడో చూపించుకున్నాను.

%%K t12-03-1 p
3. --- \VUgras{వేచి ఉండే గదిలోకి} \textsuperscript{(7)} వెళ్ళే ముందు, తిరుగు ప్రయాణంతో
కూడిన \VUgras{టికెట్టు} \textsuperscript{(9)} తీసుకున్నాను.

%%K t12-04-1 p
4. --- \VUgras{టికెట్టు కిటికీ} \textsuperscript{(8)} దగ్గర నేను ఎక్కువసేపు ఆగలేదు,
అక్కడ జనం తక్కువ. సెలవుపై బయలుదేరుతున్న \VUgras{సైనికుడు} \textsuperscript{(10)} ఒకడు
దయతో తన వంతు నాకిచ్చాడు; అందువల్ల \VUgras{స్టేషను మాస్టరును} \textsuperscript{(13)}
కొన్ని విషయాలు అడిగేంత సమయం నాకు దొరికింది — ఆయన నిఘా \VUgras{కమిషనరుతో}
\textsuperscript{(12)}, విధిలో ఉన్న \VUgras{జెండార్ముతో} \textsuperscript{(11)}
మాట్లాడుతున్నారు.

%%K t12-tit-1 sub
\VUpk{10.2pt}{40}{సామాను నమోదు.}
\VUsaut{3.0mm}

%%K t12-05-1 p
5. --- \VUgras{పుస్తకాల దుకాణంలో} \textsuperscript{(14)} పత్రిక కొని, నా
\VUgras{సామానును} \textsuperscript{(17)} తూయించాను — దాన్ని \VUgras{మోతకూలీ}
\textsuperscript{(16)} ఇప్పుడే \VUgras{సామాను బండిపై} \textsuperscript{(19)} తెచ్చాడు;
\VUgras{గుమాస్తా} \textsuperscript{(22)} — \VUgras{తూనిక యంత్రం} \textsuperscript{(21)}
ఆయన బాధ్యత — నా పెట్టె ముప్ఫై అయిదు కిలోల బరువని చెప్పాడు; అంటే అయిదు కిలోల
\VUgras{అధిక బరువు}, దానికి \VUgras{సామాను కిటికీ} \textsuperscript{(23)} దగ్గర అదనపు
రుసుము కట్టాలి.

%%K t12-06-1 p
6. --- నేను చాలా ముందుగా వచ్చాను కాబట్టి, స్టేషనులో ప్రయాణికుల రాకపోకలను చూసేంత తీరిక
దొరికింది. కొందరు తమ చిన్న సామానును \VUgras{సామాను భద్రపరిచే చోట} \textsuperscript{(25)}
పెట్టుకుంటున్నారు లేదా తిరిగి తీసుకుంటున్నారు; ఇంకొందరు \VUgras{వేళల పట్టికలు}
\textsuperscript{(26)} చూస్తున్నారు. వచ్చిన ప్రయాణికులు \VUgras{బయటికి వెళ్ళే దారికి}
\textsuperscript{(32)} హడావిడిగా వెళ్తున్నారు — చేతిలో \VUgras{సూట్‌కేసుతో}
\textsuperscript{(24)}. కొందరు \VUgras{సామాను ఇచ్చే చోట} \textsuperscript{(27)} వేచి ఉండి
తమ \VUgras{చీటీ} \textsuperscript{(29)} చూపిస్తున్నారు — తమ పెట్టెలూ, \VUgras{డబ్బాలూ}
\textsuperscript{(28)}, \VUgras{బుట్టలూ} \textsuperscript{(30)} తిరిగి తీసుకునేందుకు.

%%K t12-07-1 p
7. --- \VUgras{సుంకపు అధికారులు} \textsuperscript{(33)} మూటలను జాగ్రత్తగా పరిశీలించారు —
ప్రకటించాల్సినది ఏమైనా ఉందేమో చూసేందుకు.

%%K t12-08-1 p
8. --- కొందరు ప్రయాణికులు \VUgras{ఫలహారశాలకు} \textsuperscript{(34)} వెళ్ళారు; అక్కడ
\VUgras{సేవకుడు} \textsuperscript{(35)} వాళ్ళకు వడ్డించేందుకు హడావిడి పడుతున్నాడు. వాళ్ళు
త్వరపడాలి — \VUgras{రైలు} \textsuperscript{(36)} ఎక్స్‌ప్రెస్సు, స్టేషనులో ఎక్కువసేపు
ఆగదు.

%%K t12-09-1 p
9. --- ఆ తర్వాత బయలుదేరే ప్లాట్‌ఫారానికి వెళ్ళి, నేరుగా వెళ్ళే \VUgras{బోగీ}
\textsuperscript{(37)} కోసం వెతికాను — ప్రయాణం మధ్యలో బోగీ మార్చాల్సిన అవసరం రాకుండా.
కారిడారు బండిలోకి ఎక్కాను; అందులో పొగతాగేవాళ్ళకు ఒక \VUgras{గది} \textsuperscript{(38)},
«ఒంటరి మహిళలకు» కేటాయించిన మరో గదీ ఉన్నాయి.

%%K t12-tit-2 sub
\VUpk{10.2pt}{40}{రాత్రి బయలుదేరడం.}
\VUsaut{3.0mm}

%%K t12-10-1 p
10. --- \VUgras{మెట్టు} \textsuperscript{(40)} ఎక్కి, \VUgras{తలుపు}
\textsuperscript{(39)} మూసి, \VUgras{కిటికీ తెర} \textsuperscript{(41)} చుట్టి, నా చిన్న
సామానును \VUgras{వల అరపై} \textsuperscript{(43)} పెట్టి, \VUgras{బెంచీపై}
\textsuperscript{(42)} కూర్చున్నాను. \VUgras{దీపాలు వెలిగించేవాడు} \textsuperscript{(44)}
దీపాలు వెలిగించడం చూసి, రాత్రి బోగీలోనే గడపాల్సి వస్తుందని అనుకున్నాను;
\VUgras{తలగడల} \textsuperscript{(45)} అద్దె బాధ్యత చూసే ఉద్యోగిని పిలిచి ఒకటి అడిగాను.

%%K t12-11-1 p
11. --- టికెట్లు చూసేందుకు రైలు తనిఖీదారు వచ్చాడు. సరైన సమయం అయిపోయినా ఇంకా ఎందుకు
బయలుదేరలేదని అడిగాను. \VUgras{రైలు గార్డు} \textsuperscript{(46)} \VUgras{సామాను బోగీలో}
\textsuperscript{(47)} సైకిళ్ళు పెట్టించడంలో ఉన్నాడని చెప్పాడు. పైగా
\VUgras{కాళ్ళ వేడి పెట్టెలన్నీ} \textsuperscript{(48)} ఇంకా మార్చలేదు.

%%K t12-12-1 p
12. --- కానీ మేము త్వరలోనే బయలుదేరబోతున్నాము — \VUgras{బొగ్గు వేసేవాడు}
\textsuperscript{(50)}, తన \VUgras{బొగ్గు బండిపై} \textsuperscript{(49)} ఉండి,
\VUgras{ఇంజను} \textsuperscript{(54)} మంటను రగిల్చేందుకు బొగ్గు తీసుకున్నాడు;
\VUgras{డ్రైవరు} \textsuperscript{(51)} \VUgras{స్టేషను ఉప అధిపతి} \textsuperscript{(52)}
సంకేతం కోసమే ఎదురుచూస్తున్నాడు — ఆయన \VUgras{ప్లాట్‌ఫారంపై} \textsuperscript{(53)}
ఉన్నారు.

%%K t12-13-1 p
13. --- ఇంజను \VUgras{బాయిలరు} \textsuperscript{(56)} మెల్లగా గురగురలాడింది;
\VUgras{పొగగొట్టం} \textsuperscript{(55)} \VUgras{ఆవిరితో} \textsuperscript{(60)} కలిసిన
పొగను వరదలా కక్కింది. కీచుమనే \VUgras{కూత} \textsuperscript{(59)} మోగింది,
\VUgras{పిస్టన్లు} \textsuperscript{(58)} కదలడం మొదలుపెట్టాయి — ఆ బరువైన యంత్రపు
\VUgras{చక్రాలను} \textsuperscript{(57)} నెడుతూ.

%%K t12-tit-3 sub
\VUsaut{3.0mm}
\VUpk{10.2pt}{40}{బయలుదేరు!}
\VUsaut{3.0mm}

%%K t12-14-1 p
14. --- \VUgras{పట్టాలపై} \textsuperscript{(64)} పరుగెత్తుతున్న రైలు వేగం పెరిగింది;
\VUgras{ప్లాట్‌ఫారాలు} \textsuperscript{(65)} మాయమయ్యాయి; \VUgras{మరుగుదొడ్లను}
\textsuperscript{(66)} దాటిపోయాం, అవతలి \VUgras{పట్టాదారిపై} \textsuperscript{(63)} ఆగి
ఉన్న బోగీలనూ దాటాం : \VUgras{నిద్ర బోగీలు} \textsuperscript{(67)},
\VUgras{భోజన బోగీ} \textsuperscript{(68)}, \VUgras{తపాలా బోగీ} \textsuperscript{(69)}.

%%K t12-noto-1 noto
\VUnotes{143.2mm}{%
\textit{(*) గమనిక.} --- \textit{ఈ ప్రయాణ వర్ణన యుద్ధానికి ముందటి సరిహద్దును
బట్టి ఉందని అర్థం చేసుకోవాలి.}%
}

%%K t12-15-1 p
15. --- ఆ తర్వాత \VUgras{సరుకుల స్టేషను} \textsuperscript{(71)} మా కళ్ళ ముందునుంచి
వెళ్ళిపోయింది. పాకల కింద ఉద్యోగులు నెమ్మది రైలుతో పంపిన \VUgras{సరుకును}
\textsuperscript{(72)} ఎక్కిస్తున్నారు. \VUgras{పనివాళ్ళ జట్టు} \textsuperscript{(76)} —
\VUgras{నీటి తొట్టి} \textsuperscript{(73)} దగ్గర — \VUgras{పశువుల బోగీలనూ}
\textsuperscript{(75)}, \VUgras{చదునైన బోగీలనూ} \textsuperscript{(79)} కదిపింది,
\VUgras{సరుకు రైలు} \textsuperscript{(80)} కట్టేందుకు.

%%K t12-16-1 p
16. --- చివరి \VUgras{పట్టాల చీలికను} \textsuperscript{(78)} దాటాం, దాని దగ్గరే ముల్లు
మార్చేవాడు నిలబడి ఉన్నాడు; \VUgras{సంకేత చక్రాన్ని} \textsuperscript{(86)} దాటాం,
స్టేషను దూరంలో మాయమైంది.

%%K t12-17-1 p
17. --- త్వరలోనే \VUgras{ఇనుప వంతెన} \textsuperscript{(74)} మీదుగా వెళ్ళాం — అది
\VUgras{నదిని} \textsuperscript{(81)} దాటుతుంది. ఆ వంతెన దాటాక నది వెంట సాగాం; ఆ నదిపై
శక్తిమంతమైన \VUgras{లాగుడు నౌకలు} \textsuperscript{(82)} బరువైన \VUgras{సరుకు పడవలను}
\textsuperscript{(83)} లాగుతున్నాయి. ఇంకా \VUgras{ప్రయాణికుల ఓడను} \textsuperscript{(84)}
చూశాం — అది \VUgras{తేలే దిమ్మెకు} \textsuperscript{(85)} చేరుతున్న క్షణంలో.

%%K t12-18-1 p
18. --- చాలా స్టేషన్లను వేగంగా దాటేసి, కొన్ని \VUgras{సొరంగాల} \textsuperscript{(87)}
గుండా వెళ్ళాం; \VUgras{గేటు కాపలాదారుల} లెక్కలేనన్ని \VUgras{ఇళ్ళను}
\textsuperscript{(88)} వెనక్కి వదిలేశాం. రైలు ఊపులో ఊయలలా ఊగుతూ నేను గాఢంగా
నిద్రపోయాను.

%%K t12-tit-4 sub
\VUsaut{3.0mm}
\VUpk{10.2pt}{40}{డాయిచ్-ఆవ్రీకూర్ స్టేషను.}
\VUsaut{3.0mm}

%%K t12-19-1 p
19. --- ఉద్యోగి ఒకడు అరుస్తున్న కేకతో నాకు అకస్మాత్తుగా మెలకువ వచ్చింది :
«డాయిచ్-ఆవ్రీకూర్! \textsuperscript{(*)} అందరూ దిగండి!» అది సుంకపు తనిఖీ, సరిహద్దు
విసుగు పుట్టించే లాంఛనాలన్నీ. నా జేబు గడియారాన్ని స్టేషను \VUgras{పెద్ద గడియారానికి}
\textsuperscript{(89)} సరిచేసుకోవాల్సి వచ్చింది — అది యాభై అయిదు నిమిషాలు ముందుంది;
నిద్రమత్తు వదలని నేను \VUgras{పెద్ద ముల్లును} \textsuperscript{(90)}
\VUgras{చిన్నదానినుంచి} \textsuperscript{(91)} కష్టంగా వేరుచేసి చూశాను; ఉదయపు పొగమంచు
వల్ల \VUgras{అంకెల పలకే} \textsuperscript{(92)} పూర్తిగా అస్పష్టంగా ఉంది.

%%K t12-20-1 p
20. --- సుంకపు అధికారులు తనిఖీ చేస్తుండగా, ఆవులిస్తూ నేను స్టేషను గోడలను అలంకరించే
\VUgras{ప్రకటన పత్రాలను} \textsuperscript{(93)} చదివాను. సమయం గడిపేందుకు అల్పాహారం
తీసుకున్నాను, స్ట్రాస్‌బూర్గ్ ఇంకా ఎంత దూరమో చూసేందుకు జర్మనీ \VUgras{రైలు వల}
\textsuperscript{(94)} పటం చూశాను; నా ప్రయాణ గమ్యం త్వరలో చేరుకుంటానన్న ఆశతో బలం
పుంజుకుని మళ్ళీ నా బోగీలోకి ఎక్కాను.

\VUsaut{6.0mm}
\VUfilet{22mm}
